
\section{Aim 1 --- real-time analysis at the edge and edge-to-HPC}
\label{sec:aim1}

\textit{Strawman owners: Cappello, with Yoon (calibration, Bragg finder),
Herbst/Miceli (FPGA), Foster/Schwarz (HPC retraining). $\sim$700 words.}

Aim~1 is the inner loop plus the HPC path that the outer loop still needs.
Applications named in 2023 were ptychography, crystallography, and tomography,
with diffraction workflows alongside. Sub-deliverables included ultrafast
calibration (including a crystallography demonstration); a Bragg-peak finder
that runs on \emph{uncalibrated} images for FEL, synchrotron, and neutron
sources, with closed loop at LCLS, APS, and HFIR; transformer-based lossy
reduction and compact models on FPGAs; file-based and streaming HPC workflows;
and model re-training with a shared repository.

Edge and FPGA work does not retire HPC, and HPC does not give millisecond
steering. Both timescales remain in Fig.~\ref{fig:loops}. The section should
say what is streaming versus file-based, what is on the detector path versus
after the file, and which science metric (Bragg quality, ptycho reconstruction,
tomography artifact) the reducer actually sees.

\textit{Placeholders from the August 2026 All Hands, not an outline:} FPGA
status; compression comparisons; streaming and Tiled. LCLS crystallography
content belongs here and in Aim~3, not as the whole consortium story. Do not
quote Run~26 downtime percentages without facility operations.

\subsection{Compression}
\textit{Strawman owners: Cappello and Underwood $\sim$1200 words. for this subsection}

Why compression for light sources.

What has been explored as part of the Illumine
In the past, we have explored several compression approaches for reducing the massive data movement and storage requirements of light source experiments with Advanced Photon Source (APS). 
One major outcome is lsCOMP, a GPU-based compression framework specifically designed for the irregular unsigned integer data generated by APS applications. 
Unlike conventional compressors that are often designed for general-purpose byte streams or smooth scientific data, lsCOMP supports configurable lossless and lossy compression through adaptive scalar quantization, selective pooling, and zero-bit elimination. 
These operations are integrated into a highly optimized single-GPU-kernel design, allowing lsCOMP to achieve compression throughput of hundreds of GB/s on a single NVIDIA A100 GPU while maintaining high compression ratios and application-level data quality. 
The resulting tool has been open-sourced and published at SC 2025 (DOI: 10.1145/3712285.3759814; Code: https://github.com/szcompressor/lsCOMP).
Beyond GPU compression, we have explored how compression can be integrated more deeply into the APS end-to-end light-source data path. 
In particular, light-source workflows involve heterogeneous instruments, computing platforms, networks, and storage systems with substantially different resource and bandwidth constraints. 
This makes compression a system-level optimization problem rather than simply a standalone compression-kernel problem: an approach that performs well at one stage may not provide the best overall performance when data must traverse the complete workflow.
We have also explored learning-assisted compression to exploit redundancy that is difficult for conventional numerical compressors to capture. 
Light source datasets can exhibit complex correlations across measurements, spatial regions, and time, creating opportunities for learned prediction to reduce the amount of information that must be explicitly encoded. 
However, these correlations vary considerably across experiments and scientific applications, making it challenging to achieve high prediction accuracy, user-acceptable inference overhead, and reliable reconstruction simultaneously. 
Together, these explorations highlight the need to co-design compression algorithms with the characteristics of light source data and the systems through which those data move.


What is next for this work.

Note Lessons Learned from this

\subsection{Edge Computing}

